% English translation of chapter2.tex lines 293--518.
% Source: Methods of Algebra, Volume 2, commit
% 9a5803ff77dd3257484cb177f851a73770a59dd3 (CC BY 4.0).
% stable-unit-id: o014.aljabr2.chapter2.diagram-lemmas

% segment-id: o014.li2.chapter2.diagram-lemmas.h001
\section{Some Diagram Lemmas}\label{sec:diagram-lemmas}
\hypertarget{o014.aljabr2.chapter2.diagram-lemmas}{ }

% segment-id: o014.li2.chapter2.diagram-lemmas.p001
The arguments in this section follow \cite[Chapters 8 and 12]{KS06}. We begin
with a criterion for exactness that may be viewed as a substitute for diagram
chasing \cite[\S~6.8]{Li1} in general abelian categories. Throughout this
section, fix an abelian category $\mathcal{A}$.

% segment-id: o014.li2.chapter2.diagram-lemmas.le001
\begin{lemma}\label{prop:Abel-cat-exact-aux}
	Let $X'\xrightarrow{f}X\xrightarrow{g}X''$ be a complex. This complex is
	exact if and only if, for every morphism $S\xrightarrow{h}X$ satisfying
	$gh=0$, there are a morphism $S'\to X'$ and an epimorphism
	$S'\twoheadrightarrow S$ that make the following diagram commute:
	% segment-id: o014.li2.chapter2.diagram-lemmas.g001
	\[\begin{tikzcd}
		S' \arrow[twoheadrightarrow, r] \arrow[d] & S \arrow[d, "h"] \\
		X' \arrow[r, "f"'] & X .
	\end{tikzcd}\]
\end{lemma}

% segment-id: o014.li2.chapter2.diagram-lemmas.pf001
\begin{proof}
	First consider the ``only if'' direction. The morphisms $f$ and $h$ both
	factor through $\Ker(g)$. On the basis of these factorizations, take
	$S':=S\dtimes{\Ker(g)}X'$, and take $S'\to S$ and $S'\to X'$ to be the
	natural morphisms from the fiber product. Exactness implies that
	$X'\twoheadrightarrow\Ker(g)$, so
	Proposition~\ref{prop:Abel-cat-pull-push} implies that $S'\to S$ is also an
	epimorphism.

	Now consider the ``if'' direction. For $S:=\Ker(g)$ with its natural
	morphism $h:S\hookrightarrow X$, the hypothesis supplies the solid part of
	the following diagram:
	% segment-id: o014.li2.chapter2.diagram-lemmas.g002
	\[\begin{tikzcd}
		S' \arrow[d] \arrow[twoheadrightarrow, r] & S \arrow[d, "h"] \\
		X' \arrow[dashed, ru, "\alpha" description] \arrow[r, "f"'] & X
	\end{tikzcd}\]
	The dashed part, namely $\alpha$, comes from the factorization of $f$ given
	by $gf=0$. We show that the diagram commutes. The outer frame and the
	lower-right triangle are already known to commute. Thus the two morphisms in
	the upper-left triangle become equal after composition with $h$. Since $h$
	is a monomorphism, the upper-left triangle also commutes, proving the claim.
	Finally, the diagram shows that the composite
	$S'\to X'\xrightarrow{\alpha}S$ is an epimorphism, and hence $\alpha$ is an
	epimorphism. Thus $f$ has an epi--mono factorization
	$X'\twoheadrightarrow\Ker(g)\hookrightarrow X$, which gives
	$\Image(f)=\Ker(g)$. The complex is therefore exact.
\end{proof}

% segment-id: o014.li2.chapter2.diagram-lemmas.p002
For another approach that replaces diagram chasing in general abelian
categories, see \cite[Tag 05PL]{stacks}.

% segment-id: o014.li2.chapter2.diagram-lemmas.p003
We now discuss the Snake Lemma, which is indispensable in applications.
Consider the diagram

% segment-id: o014.li2.chapter2.diagram-lemmas.q001
\begin{equation}\label{eqn:snake}
% segment-id: o014.li2.chapter2.diagram-lemmas.g003
\begin{tikzcd}
	& \Ker' \arrow[r] \arrow[d] & \Ker \arrow[r] \arrow[d] & \Ker'' \arrow[d] & \\
	& X' \arrow[d] \arrow[r, "f"]  & X \arrow[d] \arrow[r, "g"]  & X'' \arrow[r] \arrow[d, ""{coordinate, name=Z}] & 0  \\
	0 \arrow[r] & Y' \arrow[r, "u"] \arrow[d] & Y \arrow[r, "v"] \arrow[d] & Y'' \arrow[d] & \\
	& \Coker' \arrow[uuurr, "\delta" description, crossing over, dashed, leftarrow, rounded corners, to path= {-- ([xshift=-2ex]\tikztostart.west) |- (Z) [near end]\tikztonodes -| ([xshift=2ex]\tikztotarget.east) -- (\tikztotarget) } ] \arrow[r] & \Coker \arrow[r] & \Coker'' &
\end{tikzcd}
\end{equation}

% segment-id: o014.li2.chapter2.diagram-lemmas.p004
subject to the following assumptions:
% segment-id: o014.li2.chapter2.diagram-lemmas.l001
\begin{itemize}
	% segment-id: o014.li2.chapter2.diagram-lemmas.i001
	\item $X'\xrightarrow{f}X\xrightarrow{g}X''\to0$ and
	$0\to Y'\xrightarrow{u}Y\xrightarrow{v}Y''$ are given exact sequences;
	% segment-id: o014.li2.chapter2.diagram-lemmas.i002
	\item $\Ker':=\Ker[X'\to Y']\hookrightarrow X'$, with $\Ker$ and
	$\Ker''$ defined in the same way;
	% segment-id: o014.li2.chapter2.diagram-lemmas.i003
	\item $Y'\twoheadrightarrow\Coker':=\Coker[X'\to Y']$, with $\Coker$ and
	$\Coker''$ defined in the same way;
	% segment-id: o014.li2.chapter2.diagram-lemmas.i004
	\item the morphisms $\Ker'\to\Ker$, $\Coker'\to\Coker$, and so forth come
	from the functoriality of kernels and cokernels in
	\eqref{eqn:Ker-Coker-functoriality}.
\end{itemize}

% segment-id: o014.li2.chapter2.diagram-lemmas.p005
Thus \eqref{eqn:snake} is a commutative diagram; every column and both middle
rows are exact. We now explain how to construct the dashed \emph{connecting
morphism} $\delta:\Ker''\to\Coker'$.
\index{connecting morphism}

% segment-id: o014.li2.chapter2.diagram-lemmas.p006
The first step is to derive the following commutative diagram from
\eqref{eqn:snake}; all its rows are exact:

% segment-id: o014.li2.chapter2.diagram-lemmas.q002
\begin{equation}\label{eqn:snake-conn}
% segment-id: o014.li2.chapter2.diagram-lemmas.g004
\begin{tikzcd}
	0 \arrow[r] & \Ker(g) \arrow[hookrightarrow, r, "s"] \arrow[equal, d] & X \dtimes{X''} \Ker'' \arrow[twoheadrightarrow, r] \arrow[d] \arrow[phantom, rd, "\Box" description] & \Ker'' \arrow[d] & \\
	0 \arrow[r] & \Ker(g) \arrow[hookrightarrow, r] \arrow[dashrightarrow, d] & X \arrow[r, "g"'] \arrow[d] & X'' \arrow[dashrightarrow, d] \arrow[r] & 0 \\
	0 \arrow[r] & Y' \arrow[r, "u"] \arrow[d] & Y \arrow[twoheadrightarrow, r] \arrow[d] & \Coker(u) \arrow[equal, d] \arrow[r] & 0 \\
	& \Coker' \arrow[hookrightarrow, r] & \Coker' \dsqcup{Y'} Y \arrow[r] \arrow[phantom, lu, "\boxplus" description] \arrow[twoheadrightarrow, r, "t"'] & \Coker(u) \arrow[r] & 0
\end{tikzcd}
\end{equation}

% segment-id: o014.li2.chapter2.diagram-lemmas.p007
Here the following points apply:
% segment-id: o014.li2.chapter2.diagram-lemmas.l002
\begin{compactitem}
	% segment-id: o014.li2.chapter2.diagram-lemmas.i005
	\item $\Box$ and $\boxplus$ indicate that the corresponding squares are,
	respectively, a pullback and a pushout diagram. Pullbacks preserve kernels,
	while pushouts preserve cokernels; this explains the upper-left and
	lower-right squares in the diagram;
	% segment-id: o014.li2.chapter2.diagram-lemmas.i006
	\item Proposition~\ref{prop:Abel-cat-pull-push} applies to the $\Box$ and
	$\boxplus$ parts of \eqref{eqn:snake-conn}. This explains why
	$X\dtimes{X''}\Ker''\to\Ker''$ is an epimorphism and
	$\Coker'\to\Coker'\dsqcup{Y'}Y$ is a monomorphism;
	% segment-id: o014.li2.chapter2.diagram-lemmas.i007
	\item the dashed arrow $\Ker(g)\dashrightarrow Y'$ is none other than the
	natural morphism $\Ker(g)\to\Ker(v)$ coming from the functoriality of
	kernels in \eqref{eqn:Ker-Coker-functoriality};
	% segment-id: o014.li2.chapter2.diagram-lemmas.i008
	\item similarly, the functoriality of cokernels gives the dashed arrow
	$X''\simeq\Coker(f)\to\Coker(u)$.
\end{compactitem}

% segment-id: o014.li2.chapter2.diagram-lemmas.p008
Notice that the composite
$\Ker(g)\dashrightarrow Y'\to\Coker'$ in \eqref{eqn:snake-conn} is zero,
because it becomes zero after precomposition with
$X'\twoheadrightarrow\Image(f)\rightiso\Ker(g)$. Similarly, the composite
$\Ker''\to X''\dashrightarrow\Coker(u)$ is zero, because its further
composite with $\Coker(u)\rightiso\Image(v)\hookrightarrow Y''$ is zero.

% segment-id: o014.li2.chapter2.diagram-lemmas.p009
Denote the composite along the middle path in \eqref{eqn:snake-conn} by
$\delta_0:X\dtimes{X''}\Ker''\to\Coker'\dsqcup{Y'}Y$. The preceding
observation immediately gives
% segment-id: o014.li2.chapter2.diagram-lemmas.q003
\[ \delta_0 s = 0, \quad t \delta_0 = 0. \]
Consequently, $\delta_0$ has a unique factorization
% segment-id: o014.li2.chapter2.diagram-lemmas.q004
\[ X \dtimes{X''} \Ker'' \twoheadrightarrow \Coker(s) \xrightarrow{\exists!} \Ker(t) \hookrightarrow \Coker' \dsqcup{Y'} Y . \]

The horizontal morphisms in the upper-right and lower-left corners of
\eqref{eqn:snake-conn} are known to be an epimorphism and a monomorphism,
respectively. Hence $\Coker(s)\rightiso\Ker''$ and
$\Coker'\rightiso\Ker(t)$. We have therefore obtained the desired canonical
morphism $\delta:\Ker''\to\Coker'$.

% segment-id: o014.li2.chapter2.diagram-lemmas.r001
\begin{remark}\label{rem:connecting-canonical}
	The canonical nature of the connecting morphism can be stated more
	explicitly as follows. Consider the commutative diagram
	% segment-id: o014.li2.chapter2.diagram-lemmas.g005
	\[\begin{tikzcd}[row sep=tiny, column sep=small]
		& & X' \arrow[rr] \arrow[ld] \arrow[dd] & & X \arrow[rr] \arrow[ld] \arrow[dd] & & X'' \arrow[r] \arrow[ld] \arrow[dd] & 0 \\
		0 \arrow[r] & Y' \arrow[rr, crossing over] & & Y \arrow[rr, crossing over] & & Y'' & & \\
		& & \underline{X'} \arrow[rr] \arrow[ld] & & \underline{X} \arrow[rr] \arrow[ld] & & \underline{X''} \arrow[ld] \arrow[r] & 0 \\
		0 \arrow[r] & \underline{Y'} \arrow[rr] \arrow[leftarrow, uu, crossing over] & & \underline{Y} \arrow[rr] \arrow[leftarrow, uu, crossing over] & & \underline{Y''} \arrow[leftarrow, uu, crossing over] & &
	\end{tikzcd}\]
	with every row exact. The diagram gives rise to $\Ker'$,
	$\underline{\Ker'}$, and so forth. The corresponding connecting morphisms
	$\delta$ and $\underline{\delta}$ fit into the commutative diagram
	% segment-id: o014.li2.chapter2.diagram-lemmas.g006
	\[\begin{tikzcd}
		\Ker'' \arrow[r, "\delta"] \arrow[d, "\text{natural morphism}"'] & \Coker' \arrow[d, "\text{natural morphism}"] \\
		\underline{\Ker''} \arrow[r, "\underline{\delta}"'] & \underline{\Coker'} .
	\end{tikzcd}\]
	By definition, to prove this it suffices to show that $\delta_0$ and
	$\underline{\delta_0}$ also fit into the corresponding commutative square.
	Everything follows from the functoriality of kernels, cokernels, products,
	and coproducts. The details are left to the reader.
\end{remark}

% segment-id: o014.li2.chapter2.diagram-lemmas.p010
The following two observations will be used in the proof below.
% segment-id: o014.li2.chapter2.diagram-lemmas.l003
\begin{itemize}
	% segment-id: o014.li2.chapter2.diagram-lemmas.i009
	\item The construction of the connecting morphism is self-dual: when
	\eqref{eqn:snake} is viewed in $\mathcal{A}^{\opp}$, all arrows are
	reversed, while the roles of $\Ker'$ and $\Coker'$, and so forth, are
	interchanged. The resulting diagram in $\mathcal{A}^{\opp}$ still resembles
	\eqref{eqn:snake}, except that it is rotated through half a turn. When the
	corresponding connecting morphism is viewed back in $\mathcal{A}$, it still
	runs from $\Ker''$ to $\Coker'$, and it agrees with the $\delta$ constructed
	above.
	
	% segment-id: o014.li2.chapter2.diagram-lemmas.i010
	\item The morphism $\Ker\to\Ker''$ factors as
	$\Ker\to X\dtimes{X''}\Ker''\to\Ker''$. Inspection of
	\eqref{eqn:snake-conn} shows that the two composites
	% segment-id: o014.li2.chapter2.diagram-lemmas.q005
	\[ \Ker \to \Ker'' \xrightarrow{\delta} \Coker' \hookrightarrow \Coker' \dsqcup{Y'} Y, \qquad \Ker \to X \dtimes{X''} \Ker'' \xrightarrow{\delta_0} \Coker' \dsqcup{Y'} Y \]
		are equal. But the second composite is also equal to
	$\Ker\to X\to Y\to\Coker'\dsqcup{Y'}Y$, which is zero. We therefore obtain
	% segment-id: o014.li2.chapter2.diagram-lemmas.q006
	\begin{equation}\label{eqn:snake-conn-complex}
		\Ker \to \Ker'' \xrightarrow{\delta} \Coker' \quad \text{has composite}\; 0.
	\end{equation}
\end{itemize}

% segment-id: o014.li2.chapter2.diagram-lemmas.th001
\begin{theorem}[Snake Lemma]\label{prop:snake-lemma}\index{Snake Lemma}
	Consider diagram \eqref{eqn:snake} in an abelian category $\mathcal{A}$.
	Then
	$\Ker'\to\Ker\to\Ker''\xrightarrow{\delta}\Coker'\to\Coker\to\Coker''$
	is an exact sequence. If $f:X'\to X$ in the diagram is a monomorphism (or
	$v:Y\to Y''$ is an epimorphism), then $\Ker'\to\Ker$ is also a
	monomorphism (or $\Coker\to\Coker''$ is also an epimorphism).
\end{theorem}
% segment-id: o014.li2.chapter2.diagram-lemmas.pf002
\begin{proof}
	The last assertion is easy to prove: if $f$ is a monomorphism, then the
	composite $\Ker'\hookrightarrow X'\xrightarrow{f}X$ is also a
	monomorphism. Since \eqref{eqn:snake} commutes, this implies that
	$\Ker'\to\Ker$ is a monomorphism. The case in which $v$ is an epimorphism
	is obtained by reversing all arrows.

	The main part of the proof is the exactness of the first sequence. By the
	duality observed above, it suffices to prove that
	$\Ker'\to\Ker\to\Ker''\to\Coker'$ is exact.
	
	First, $\Ker'\to\Ker\to\Ker''$ is a complex. Indeed, it is induced by
	$X'\to X\to X''$, whose composite is zero. We now use
	Lemma~\ref{prop:Abel-cat-exact-aux} to prove exactness. Suppose that a
	morphism $\psi:S\to\Ker$ makes the composite
	$S\xrightarrow{\psi}\Ker\to\Ker''$ zero. We must construct a commutative
	diagram
	% segment-id: o014.li2.chapter2.diagram-lemmas.q007
	\begin{equation}\label{eqn:snake-lemma-aux-1}
	% segment-id: o014.li2.chapter2.diagram-lemmas.g007
	\begin{tikzcd}
		S' \arrow[twoheadrightarrow, r, "h"] \arrow[d] & S \arrow[d, "\psi"] \arrow[rd, "0"] & \\
		\Ker' \arrow[r] \arrow[hookrightarrow, d] & \Ker \arrow[hookrightarrow, d] \arrow[r] & \Ker'' \arrow[hookrightarrow, d] \\
		X' \arrow[r, "f"'] & X \arrow[r, "g"'] & X''
	\end{tikzcd}\end{equation}
	where $\Ker'\leftarrow S'\xrightarrow{h}S$ remains to be constructed. To
	do so, first apply Lemma~\ref{prop:Abel-cat-exact-aux} to the composite
	$S\xrightarrow{\psi}\Ker\to X$ to obtain the commutative diagram
	% segment-id: o014.li2.chapter2.diagram-lemmas.g008
	\[\begin{tikzcd}
		S' \arrow[twoheadrightarrow, r, "h"] \arrow[d] & S \arrow[d] \arrow[rd, "0"] & \\
		X' \arrow[r, "f"'] & X \arrow[r, "g"'] & X'' .
	\end{tikzcd}\]
	Thus the composite $S'\to X'\to Y'\xrightarrow{u}Y$ is equal to the
	composite $S'\xrightarrow{\psi h}\Ker\to X\to Y$, namely zero. Since $u$
	is a monomorphism, the composite $S'\to X'\to Y'$ is also zero. Hence
	$S'\to X'$ factors through $\Ker'$, giving all the arrows in
	\eqref{eqn:snake-lemma-aux-1}.
	
	Now consider the left part of \eqref{eqn:snake-lemma-aux-1}: the outer
	frame and the lower square commute. The familiar
	Lemma~\ref{prop:epi-mono-morphism} shows that the upper square also
	commutes. Therefore $\Ker'\to\Ker\to\Ker''$ is exact.
	
	Next consider $\Ker\to\Ker''\xrightarrow{\delta}\Coker'$. Equation
	\eqref{eqn:snake-conn-complex} first shows that this sequence is a complex.
	To prove its exactness, we again use Lemma~\ref{prop:Abel-cat-exact-aux}.
	Suppose that a morphism $\psi:S\to\Ker''$ satisfies $\delta\psi=0$. We
	must construct a commutative diagram of the form
	% segment-id: o014.li2.chapter2.diagram-lemmas.q008
	\begin{equation}\label{eqn:snake-lemma-aux-2}
	% segment-id: o014.li2.chapter2.diagram-lemmas.g009
	\begin{tikzcd}
		S_0 \arrow[d] \arrow[twoheadrightarrow, r] & S \arrow[d, "\psi"] \\
		\Ker \arrow[r] & \Ker''
	\end{tikzcd}\end{equation}
	The construction of $\delta$ showed that
	$X\dtimes{X''}\Ker''\to\Ker''\to0$ is exact. Hence there is a commutative
	diagram
	% segment-id: o014.li2.chapter2.diagram-lemmas.q009
	\begin{equation}\label{eqn:snake-lemma-aux-3}
	% segment-id: o014.li2.chapter2.diagram-lemmas.g010
	\begin{tikzcd}
		S_1 \arrow[twoheadrightarrow, r] \arrow[d] & S \arrow[d, "\psi"'] \arrow[rd, "0"] & \\
		X \dtimes{X''} \Ker'' \arrow[r] & \Ker'' \arrow[r] & 0 .
	\end{tikzcd}\end{equation}

	Consider the composite
	$S_1\to X\dtimes{X''}\Ker''\to X\to Y\xrightarrow{v}Y''$. Comparing the
	diagram above with \eqref{eqn:snake} and \eqref{eqn:snake-conn}, we see that
	this composite is zero. Hence $S_1\to X\dtimes{X''}\Ker''\to X\to Y$
	factors as $S_1\xrightarrow{\xi}Y'\xrightarrow{u}Y$. We now prove that the
	composite $S_1\xrightarrow{\xi}Y'\to\Coker'$ is zero. Begin with the
	commutative diagram
	% segment-id: o014.li2.chapter2.diagram-lemmas.g011
	\[\begin{tikzcd}
		S \arrow[r, "\psi"] & \Ker'' \arrow[r, "\delta"] & \Coker' \arrow[hookrightarrow, d] \\
		S_1 \arrow[twoheadrightarrow, u] \arrow[r] & X \dtimes{X''} \Ker'' \arrow[r, "\delta_0"] \arrow[u] \arrow[d] & \Coker' \dsqcup{Y'} Y \\
		& X \arrow[r] & Y \arrow[u]
	\end{tikzcd}\]

	The composite along the first row is $\delta\psi=0$, so the composite along
	the second row is also zero. Consequently, the composite
	$S_1\xrightarrow{\xi}Y'\xrightarrow{u}Y\to\Coker'\dsqcup{Y'}Y$ is zero.
	On the other hand, the composite
	$Y'\xrightarrow{u}Y\to\Coker'\dsqcup{Y'}Y$ equals
	$Y'\to\Coker'\to\Coker'\dsqcup{Y'}Y$. In the construction of $\delta$ we
	proved that $\Coker'\to\Coker'\dsqcup{Y'}Y$ is a monomorphism. Thus the
	composite $S_1\xrightarrow{\xi}Y'\to\Coker'$ is indeed zero.

	The next step is to apply Lemma~\ref{prop:Abel-cat-exact-aux} to the exact
	sequence $X'\to Y'\to\Coker'$. This gives a commutative diagram
	% segment-id: o014.li2.chapter2.diagram-lemmas.g012
	\[\begin{tikzcd}
		S_0 \arrow[twoheadrightarrow, r] \arrow[d, "k"'] & S_1 \arrow[d, "\xi"] \arrow[rd, "0"] & \\
		X' \arrow[r] & Y' \arrow[r] & \Coker' .
	\end{tikzcd}\]

	Write $\lambda$ for the composite
	$S_0\twoheadrightarrow S_1\to X\dtimes{X''}\Ker''\to X$. Also write the
	morphisms $X'\to Y'$ and $X\to Y$ as $a$ and $b$, respectively. The
	preceding discussion gives the commutative diagram
	% segment-id: o014.li2.chapter2.diagram-lemmas.g013
	\[\begin{tikzcd}
		S_0 \arrow[twoheadrightarrow, r] \arrow[d, "k"' near end] \arrow[rrd, "\lambda" description, near end] & S_1 \arrow[d, crossing over, "\xi"' near end] \arrow[r] & X \dtimes{X''} \Ker'' \arrow[d] \\
		X' \arrow[r, "a"] \arrow[d, "f"'] & Y' \arrow[d, "u"'] & X \arrow[ld, "b"] \\
		X \arrow[r, "b"] & Y &
	\end{tikzcd}\]
	Both squares in this diagram are known to commute; commutativity of the
	triangle comes from the definition of $\lambda$, while commutativity of the
	trapezoid on the right comes from the preceding discussion of $\xi$.
	
	It follows at once that $b\lambda=bfk$, so
	$\lambda-fk:S_0\to X$ factors through $\Ker=\Ker(b)\hookrightarrow X$.
	Finally, consider the diagram
	% segment-id: o014.li2.chapter2.diagram-lemmas.g014
	\[\begin{tikzcd}
		S_0 \arrow[twoheadrightarrow, r] \arrow[d] \arrow[dd, bend right=50, "{\lambda -fk}"'] & S_1 \arrow[twoheadrightarrow, r] & S \arrow[d, "\psi"] \\
		\Ker \arrow[hookrightarrow, d] \arrow[rr] & & \Ker'' \arrow[hookrightarrow, d] \\
		X \arrow[rr, "g"] & & X'' .
	\end{tikzcd}\]
	To obtain the diagram in \eqref{eqn:snake-lemma-aux-2}, it suffices to show
	that its upper square commutes. The lower square and the curved part are
	known to commute; the usual argument reduces the problem to commutativity of
	the outer frame. Since $g(\lambda-fk)=g\lambda$, this in turn reduces to the
	known commutative diagram (see \eqref{eqn:snake-lemma-aux-3}):
	% segment-id: o014.li2.chapter2.diagram-lemmas.g015
	\[\begin{tikzcd}
		S_0 \arrow[twoheadrightarrow, r] \arrow[rdd, bend right, "\lambda"'] & S_1 \arrow[twoheadrightarrow, r] \arrow[d] & S \arrow[d, "\psi"] \\
		& X \dtimes{X''} \Ker'' \arrow[r] \arrow[d] & \Ker'' \arrow[d] \\
		& X \arrow[r, "g"] & X'' .
	\end{tikzcd}\]
	All the required conditions are therefore satisfied.
\end{proof}

% segment-id: o014.li2.chapter2.diagram-lemmas.p011
The connecting homomorphism $\delta$ and Theorem~\ref{prop:snake-lemma} are
much simpler in the category $R\dcate{Mod}$ of $R$-modules; see
\cite[Proposition 6.8.6]{Li1}.

% segment-id: o014.li2.chapter2.diagram-lemmas.pr001
\begin{proposition}[Five Lemma]\label{prop:5-lemma}
	\index{Five Lemma}
	Consider a commutative diagram with exact rows in an abelian category
	$\mathcal{A}$:
	% segment-id: o014.li2.chapter2.diagram-lemmas.g016
	\[\begin{tikzcd}
		X_1 \arrow[r] \arrow[d, "f_1"'] & X_2 \arrow[r] \arrow[d, "f_2"'] & X_3 \arrow[r] \arrow[d, "f_3" description] & X_4 \arrow[r] \arrow[d, "f_4"] & X_5 \arrow[d, "f_5"] \\
		Y_1 \arrow[r] & Y_2 \arrow[r] & Y_3 \arrow[r] & Y_4 \arrow[r] & Y_5
	\end{tikzcd}\]
	% segment-id: o014.li2.chapter2.diagram-lemmas.l004
	\begin{enumerate}[(i)]
		% segment-id: o014.li2.chapter2.diagram-lemmas.i011
		\item If $f_1$ is an epimorphism and $f_2,f_4$ are monomorphisms, then
		$f_3$ is a monomorphism (only the first four columns are involved);
		% segment-id: o014.li2.chapter2.diagram-lemmas.i012
		\item if $f_5$ is a monomorphism and $f_2,f_4$ are epimorphisms, then
		$f_3$ is an epimorphism (only the last four columns are involved);
		% segment-id: o014.li2.chapter2.diagram-lemmas.i013
		\item if $f_1$ is an epimorphism, $f_5$ is a monomorphism, and $f_2,f_4$
		are both isomorphisms, then $f_3$ is an isomorphism.
	\end{enumerate}
\end{proposition}
% segment-id: o014.li2.chapter2.diagram-lemmas.pf003
\begin{proof}
	Clearly (i) and (ii) are dual to each other, while (iii) follows from
	(i)--(ii) and Proposition~\ref{prop:strict-isomorphism}. It therefore
	suffices to prove (i).
	
	Suppose that a morphism $h:S\to X_3$ satisfies $f_3h=0$; we must prove
	that $h=0$. The composite
	$S\xrightarrow{h}X_3\to X_4\xrightarrow{f_4}Y_4$ is zero, and therefore so
	is $S\xrightarrow{h}X_3\to X_4$. Applying
	Lemma~\ref{prop:Abel-cat-exact-aux} to the exact sequence
	$X_2\to X_3\to X_4$ gives a commutative diagram
	% segment-id: o014.li2.chapter2.diagram-lemmas.g017
	\[\begin{tikzcd}
		S' \arrow[twoheadrightarrow, r] \arrow[d] & S \arrow[d, "h"] \\
		X_2 \arrow[r] & X_3 .
	\end{tikzcd}\]

	We now construct a commutative diagram
	% segment-id: o014.li2.chapter2.diagram-lemmas.g018
	\[\begin{tikzcd}[row sep=small]
		S''' \arrow[twoheadrightarrow, r] \arrow[dd] & S'' \arrow[twoheadrightarrow, r] \arrow[dd] & S' \arrow[d] \\
		& & X_2 \arrow[d, "f_2" inner sep=0.5em] \\
		X_1 \arrow[r, "f_1"'] & Y_1 \arrow[r] & Y_2 . 
	\end{tikzcd}\]
	Here is how. Since $f_3h=0$, the composite
	$S'\to X_2\xrightarrow{f_2}Y_2\to Y_3$ is zero. Applying
	Lemma~\ref{prop:Abel-cat-exact-aux} to $S'\to Y_2$ and the exact sequence
	$Y_1\to Y_2\to Y_3$ gives the right-hand part of the diagram. Applying
	Lemma~\ref{prop:Abel-cat-exact-aux} to $S''\to Y_1$ and the exact sequence
	$X_1\xrightarrow{f_1}Y_1\to0$ gives the left-hand part.
	
	Take the composite $S'''\to S'$ and consider the following diagram:
	% segment-id: o014.li2.chapter2.diagram-lemmas.g019
	\[\begin{tikzcd}
		S''' \arrow[twoheadrightarrow, r] \arrow[d] & S' \arrow[twoheadrightarrow, r] \arrow[d] & S \arrow[d, "h"] & \\
		X_1 \arrow[r] \arrow[d, "f_1"'] & X_2 \arrow[r] \arrow[d, "f_2"] & X_3 \arrow[r] & X_4 \\
		Y_1 \arrow[r] & Y_2 & &
	\end{tikzcd}\]
	We show that this diagram commutes. The only issue is the upper-left square.
	Compose the two paths
	% segment-id: o014.li2.chapter2.diagram-lemmas.g020
	\begin{tikzpicture}[scale=0.3]
		\draw[->] (0,1) -- (1,1) -- (1,0); \end{tikzpicture}
	and
	% segment-id: o014.li2.chapter2.diagram-lemmas.g021
	\begin{tikzpicture}[scale=0.3]
		\draw[->] (0,1) -- (0,0) -- (1,0);
	\end{tikzpicture}
	in that square with $f_2$ on the right. Since the two squares
	% segment-id: o014.li2.chapter2.diagram-lemmas.g022
	$\begin{tikzcd}[row sep=small, column sep=small]
		X_1 \arrow[r] \arrow[d] & X_2 \arrow[d] \\
		Y_1 \arrow[r] & Y_2
	\end{tikzcd}$ and
	% segment-id: o014.li2.chapter2.diagram-lemmas.g023
	$\begin{tikzcd}[row sep=small, column sep=small]
		S''' \arrow[r] \arrow[d] & S' \arrow[d] \\
		Y_1 \arrow[r] & Y_2
	\end{tikzcd}$
	commute and $f_2$ is a monomorphism, the upper-left square also commutes.
	
	Thus the composite
	$S'''\twoheadrightarrow S'\twoheadrightarrow S\xrightarrow{h}X_3$ is zero,
	and hence $h=0$. This proves the result.
\end{proof}

% segment-id: o014.li2.chapter2.diagram-lemmas.p012
This section has introduced only two of the most commonly used diagram lemmas.
Using the language of double complexes, G.\ Bergman found the comprehensive
Salamander Lemma, which subsumes the Snake Lemma and several classical diagram
lemmas of homological algebra. Interested readers may consult \cite{Be12}.
\index{Salamander Lemma}
